% Katalog E: Inline-Semantik und Verweise (Palette Slot 5).
\ZSFNewColumn
\StartChapter[kat:marker]{E · Marker}

\begin{runintext}
  Auszeichnungen im laufenden Text. Jeder Eintrag trägt denselben
  Trägersatz, damit sich die Marker direkt vergleichen lassen — vier von
  ihnen sehen absichtlich gleich aus und unterscheiden sich nur in der
  Bedeutung.
\end{runintext}

\SubsectionBar[kat:marker-betonung]{Betonung}
\ZSFindex{ZSFkeyword}
\ZSFindex{ZSFlabel}

\begin{defbox}[E-01 · ZSFkeyword][weight=quiet]
  Mustertext mit \ZSFkeyword{Fachbegriff} zum Formvergleich.
\end{defbox}
\Zweck{Zentraler Fachbegriff als Scan-Anker; erzeugt automatisch einen
Registereintrag. Richtwert 1--3 pro Block.}

\begin{defbox}[E-02 · ZSFkeyword*][weight=quiet]
  Mustertext mit \ZSFkeyword*{Fachbegriff} zum Formvergleich.
\end{defbox}
\Zweck{Dasselbe ohne Registereintrag — für Begriffe, die nur hier stehen.}

\begin{defbox}[E-03 · ZSFlabel][weight=quiet]
  Mustertext mit \ZSFlabel{Beschriftung} zum Formvergleich.
\end{defbox}
\Zweck{Neutrale fette Beschriftung ohne Fachbegriff-Semantik: Listenkopf,
Zeilenmarke. Sieht aus wie E-01 und ist es nicht.}

\begin{defbox}[E-04 · ZSFhl][weight=quiet]
  Mustertext mit \ZSFhl{hervorgehobener Kernaussage} zum Formvergleich.
\end{defbox}
\Zweck{Die prüfungskritischste Aussage einer Box. Sparsam, sonst flacht das
Signal ab.}

\begin{defbox}[E-05 · ZSFdanger][weight=quiet]
  Mustertext mit \ZSFdanger{Stolperfalle} zum Formvergleich.
\end{defbox}
\Zweck{Inline-Pille für eine kurze Warnung innerhalb einer anderen Box —
nicht zusätzlich zur warnbox für dieselbe Aussage.}

\begin{defbox}[E-06 · ZSFconclusion][weight=quiet]
  Mustertext zum Formvergleich. \ZSFconclusion{Daraus folgt die Aussage.}
\end{defbox}
\Zweck{Leitet eine Folgerung ein, gesetzt als Implikationspfeil.}

\SubsectionBar[kat:marker-verweis]{Verweise}
\ZSFindex{ZSFref}
\ZSFindex{ZSFsectionref}

\begin{defbox}[E-07 · ZSFref][weight=quiet]
  Mustertext mit Verweis \ZSFref{kat:boxen} zum Formvergleich.
\end{defbox}
\Zweck{Querverweis mit Pfeil, in der Farbe des Zielkapitels — der einzige
farbtragende Marker im Fliesstext.}

\begin{defbox}[E-08 · ZSFsectionref][weight=quiet]
  Mustertext mit Zielnummer \ZSFsectionref{kat:tabellen} zum Formvergleich.
\end{defbox}
\Zweck{Kompakte klickbare Zielnummer ohne Pfeil — für lokale
Inhaltsübersichten und Referenztabellen.}

\begin{defbox}[E-09 · ZSFScriptRef][weight=quiet]
  Mustertext zum Formvergleich.
  \ZSFScriptRef{42}
\end{defbox}
\Zweck{Rechtsbündiger Verweis auf die Skript-Seite; nur wo das Skript in der
Prüfung erlaubt ist.}

\begin{defbox}[E-10 · ZSFTitleTag\ZSFTitleTag{hier}][weight=quiet]
  Mustertext zum Formvergleich — das Tag steht rechts im Titel oben.
\end{defbox}
\Zweck{Rechtsbündiges Meta-Tag im Titel jeder Box mit Titel, auch in der
leisen Fassung.}

\begin{defbox}[E-11 · Ink-Vertrag: \ZSFref{kat:formeln} im Titel]
  Derselbe Verweis im Inhalt: \ZSFref{kat:formeln} — hier mit seiner Farbe.
\end{defbox}
\Zweck{Auf einer Fläche mit eigener Kontrastfarbe geben farbtragende Marker
ihre Farbe ab. Kein Regler, sondern automatisch.}

\SubsectionBar[kat:marker-block]{Blockgliederung}
\ZSFindex{ZSFsep}
\ZSFindex{ZSFItem}

\begin{defbox}[E-12 · ZSFsep][weight=quiet]
  Mustertext zum Formvergleich.
  \ZSFsep
  Zweiter Block nach dem Trenner.
  \ZSFsep[E-13 · ZSFsep mit Beschriftung]
  Dritter Block, mit benanntem Trenner darüber.
\end{defbox}
\Zweck{Trennt Blöcke innerhalb jeder Box — nicht nur in der formulabox. Wer
stattdessen eine zweite Box daneben stellt, hat den Trenner nicht gefunden.}

\begin{ZSFlist}[E-14 · ZSFItem]
  \ZSFItem{Mit Marke}{Mustertext zum Formvergleich.}
  \ZSFItem{Mustertext ohne Marke — keine Marke, kein Trenner.}
\end{ZSFlist}
\Zweck{Listeneintrag mit optionaler Marke; eine Marke zu erfinden, damit die
Form aufgeht, wäre eine Inhaltsänderung.}

\textVorBox{E-15 · textVorBox — dieser Satz gehört zur Box darunter.}
\begin{defbox}[][weight=quiet]
  Mustertext zum Formvergleich.
\end{defbox}
\Zweck{Bindet einen Satz an die folgende Box; textNachBox an die
vorhergehende — jede Zweckzeile dieses Katalogs ist eine.}

\SubsectionBar[kat:marker-farbe]{Farbe als Bedeutung}
\ZSFindex{Grössenfarbe}
\ZSFindex{ZSFfontDiagramLabel}

\begin{formulabox}[E-16 · Grössenfarben]
  \[
    \ZSFqMoment{M} = \ZSFqKraft{F} \cdot \ZSFqWeg{r}
  \]
  \formulanote{Mustertext zum Formvergleich.}
\end{formulabox}
\Zweck{Eine Farbe gehört im ganzen Dokument einer Grösse. Zwölf Slots,
vergeben in preamble.tex, nie im Kapitel.}

\begin{figbox}[E-17 · ZSFfontDiagramLabel]
  {\ZSFfontDiagramLabel Mustertext in der Beschriftungsgrösse.}\\
  {\ZSFfontDiagramLabelSmall Zweite Stufe für gedrängte Zusatzangaben.}
\end{figbox}
\Zweck{Die zwei Schriftstufen für Diagramm-Beschriftung statt lokalem tiny;
beide zusammen skalierbar über ZSFDiagramLabelScale.}
